\documentclass[11pt]{article}
\usepackage[margin=1in]{geometry}
\usepackage{amsmath, amssymb}
\usepackage{booktabs}
\usepackage{parskip}
\usepackage{setspace}
\usepackage{xcolor}
\usepackage{hyperref}

\onehalfspacing

\title{\textbf{Diaspora-Adjusted Importation Risk:\\
       Methods, Logic, and Interpretation}}
\author{Technical note --- FIFA World Cup 2026 importation risk analysis}
\date{}

\begin{document}
\maketitle
\thispagestyle{empty}

% ============================================================
\section*{1. What the core model gives us — and what it misses}
% ============================================================

The core Poisson model produces, for each source country $c$, venue city $v$,
and disease $d$:
\[
  \lambda_{c,v,d} = N_{c,v} \cdot \rho_d \cdot \frac{C_{c,d}}{P_c} \cdot p_d
\]
where $N_{c,v}$ is the projected number of travelers from $c$ to $v$,
$\rho_d$ is the under-reporting correction, $C_{c,d}/P_c$ is reported
per-capita incidence, and $p_d$ is the probability of traveling while
infectious. Summing across source countries gives the city-level
importation intensity $\Lambda_{v,d} = \sum_c \lambda_{c,v,d}$, and
the probability of at least one importation follows from the Poisson CDF:
\[
  P(\geq 1)_{v,d} = 1 - e^{-\Lambda_{v,d}}
\]

This answers the question: \textit{how many cases will arrive, and how
likely is at least one?}

What it does \textbf{not} answer is: \textit{what social environment
do arriving cases enter?} An infectious traveler from Brazil who lands in
Boston enters a city where 7.9\% of all foreign-born residents are
Brazilian --- a large, socially connected community of co-nationals who
share language, neighborhoods, community events, and World Cup watch
gatherings. The same traveler landing in Kansas City enters a city with
almost no such community. The Poisson intensity $\lambda$ is identical
in both cases if the travel volumes are the same; the downstream
transmission potential is not.

The diaspora extension quantifies this difference.

% ============================================================
\section*{2. Diaspora data and notation}
% ============================================================

We use the US Census Bureau American Community Survey (ACS) 5-year
estimates (2019--2023), Table B05006 (\textit{Place of Birth for the
Foreign-Born Population}), for the Core-Based Statistical Area (CBSA)
corresponding to each of the 11 WC host metro areas.

\begin{itemize}
  \item $D_{c,v}$ --- number of US residents born in country $c$ currently
    living in metro area $v$ (ACS estimate).
  \item $\mathrm{FB}_{v}$ --- total foreign-born population of metro $v$
    (the row total of Table B05006). This is the denominator: it represents
    \emph{all} immigrants in that city, regardless of origin country.
  \item $\kappa_{c,v} = \dfrac{D_{c,v}}{\mathrm{FB}_{v}}$ ---
    \textbf{diaspora concentration}: the share of metro $v$'s
    foreign-born population that was born in country $c$. This is a
    number between 0 and 1.
\end{itemize}

\medskip
\noindent\textbf{Concrete examples.}
Brazilians constitute 7.9\% of Boston's foreign-born population
($\kappa_{\text{Brazil, Boston}} = 0.079$). Haitians constitute 8.8\%
of Miami's ($\kappa_{\text{Haiti, Miami}} = 0.088$). Ecuadorians
constitute 4.8\% of New York's ($\kappa_{\text{Ecuador, NY}} = 0.048$).
For most country--city pairs $\kappa$ is much smaller, often below 1\%.

A critical point: $\kappa$ is normalized by the \emph{total foreign-born}
population, not by the total metro population. This makes $\kappa$ a
measure of how concentrated the source-country community is
\emph{within the immigrant population} of that city --- the relevant
social network for co-national mixing.

% ============================================================
\section*{3. Mechanism A: local social mixing at the venue city}
% ============================================================

\subsection*{3.1 Reasoning}

When an infectious traveler from country $c$ arrives at venue city $v$,
they do not mix at random with the entire population. They are more
likely to interact with co-nationals --- through shared language, cultural
spaces, restaurants, community events, and World Cup gatherings organized
within the diaspora. We model this selective mixing by assuming that, upon
arrival, each traveler's social contacts are drawn proportionally from the
foreign-born population of the city.

\subsection*{3.2 The mixing assumption}

\textbf{Assumption:} each infectious traveler from country $c$ who arrives
at city $v$ interacts with the foreign-born population of $v$. Under
proportional mixing, the probability that a randomly chosen interaction
partner was born in country $c$ (i.e., is a co-national) equals
$\kappa_{c,v}$.

This is a simplification --- actual mixing rates depend on neighborhood
clustering, social network structure, and event attendance --- but it
provides a transparent, data-grounded lower bound on co-national contact
probability.

\subsection*{3.3 Mathematics: Poisson thinning}

Let $X_{c,v,d}$ be the number of infectious travelers from country $c$
arriving at city $v$. From the core model:
\[
  X_{c,v,d} \sim \text{Poisson}(\lambda_{c,v,d})
\]
Each of these $X_{c,v,d}$ travelers independently enters the co-national
diaspora network with probability $\kappa_{c,v}$ (the mixing assumption
above). Let $Y_{c,v,d}$ be the number of those travelers who actually
enter the co-national network. By the \textbf{Poisson thinning theorem}:
\[
  Y_{c,v,d} \sim \text{Poisson}(\lambda_{c,v,d} \cdot \kappa_{c,v})
\]
We define:
\begin{equation}
  \Omega^A_{c,v,d} = \lambda_{c,v,d} \cdot \kappa_{c,v}
    = \lambda_{c,v,d} \cdot \frac{D_{c,v}}{\mathrm{FB}_{v}}
  \label{eq:omegaA}
\end{equation}

$\Omega^A_{c,v,d}$ is the expected number of imported cases of disease
$d$, from country $c$, that enter the co-national social network in
city $v$. It is a valid Poisson rate for a \emph{sub-population} of
the original $\lambda_{c,v,d}$ arrivals.

The city-level aggregate across all source countries is:
\[
  \Omega^A_{v,d} = \sum_c \Omega^A_{c,v,d}
    = \sum_c \lambda_{c,v,d} \cdot \kappa_{c,v}
\]

Because this is a sum of independent Poisson random variables,
$\Omega^A_{v,d}$ is itself a Poisson rate, and the probability that at
least one imported case enters \emph{any} co-national network in city $v$
is:
\begin{equation}
  P^A(\geq 1)_{v,d} = 1 - e^{-\Omega^A_{v,d}}
  \label{eq:probA}
\end{equation}

\subsection*{3.4 How to read and compare the two probabilities}

Figure~1 shows $P(\geq 1)_{v,d}$ --- the probability that at least one
case arrives at city $v$. Figure~4 shows $P^A(\geq 1)_{v,d}$ --- the
probability that at least one of those cases enters a co-national network.
Both are derived from the same Poisson logic. Their difference is
interpretable:

\medskip
\noindent\fbox{\parbox{0.93\textwidth}{%
\textbf{Example.} Dengue at Miami: $P(\geq 1) \approx 1$ (virtually
certain that at least one dengue case arrives). Suppose $P^A(\geq 1) =
0.60$. This means: \emph{it is virtually certain that cases arrive, but
there is only a 60\% probability that at least one of those cases enters
a co-national social network}. In 40\% of scenarios, cases arrive but
disperse into the general population --- still a public health concern,
but without the added amplification of concentrated co-national social
ties.
}}

\medskip
The gap between $P(\geq 1)$ and $P^A(\geq 1)$ is the fraction of risk
that is \emph{diffuse} (general population) rather than
\emph{concentrated} (co-national network). Cities where this gap is
small have highly concentrated diaspora communities relative to total
importation intensity; cities where the gap is large have high importation
volumes but dispersed or thin co-national networks.

\textbf{Why are $P^A$ values always lower than $P$?} Because
$\Omega^A_{v,d} \leq \Lambda_{v,d}$ always: $\kappa_{c,v} \leq 1$
for all $c$, $v$, so multiplying $\lambda_{c,v,d}$ by $\kappa_{c,v}$
can only reduce it. The diaspora network is always a sub-population of
the full arrival stream; it cannot have more cases than the whole.

% ============================================================
\section*{4. Mechanism B: diaspora return seeding}
% ============================================================

\subsection*{4.1 Reasoning}

Mechanism A asks what happens when international travelers arrive at venue
cities where co-nationals already live. Mechanism B runs in the opposite
direction: US-resident diaspora members --- already living in the US ---
travel \emph{domestically} to a WC venue city to watch their country's
matches, are exposed to infectious travelers from their home country at
the venue, and then return to their home city. The risk lands not at
the venue, but in the diaspora hub city.

This mechanism is entirely invisible to the core model, which only tracks
international arrivals. A city that hosts no WC matches and receives few
international visitors from a given country can still face secondary
seeding through returning diaspora members.

\subsection*{4.2 Notation}

Let $v_{\text{match}}$ be the venue city where country $c$ plays, and
$v_{\text{home}}$ be a hub city where a diaspora community from $c$ lives
($v_{\text{home}} \neq v_{\text{match}}$).

\subsection*{4.3 Mathematics}

The exposure source at $v_{\text{match}}$ is $\lambda_{c,\,v_{\text{match}},\,d}$
--- the expected number of infectious travelers from country $c$ arriving
at the match venue. This is the same Poisson rate computed by the
schedule-driven model (M3). Diaspora members from $v_{\text{home}}$ who
attend the match are exposed to this infectious pool. Upon return to
$v_{\text{home}}$, they may seed transmission among co-nationals there.
The secondary seeding potential is:
\begin{equation}
  \Omega^B_{c,\,v_{\text{match}},\,v_{\text{home}},\,d}
    = \lambda_{c,\,v_{\text{match}},\,d}
      \cdot \kappa_{c,\,v_{\text{home}}}
    = \lambda_{c,\,v_{\text{match}},\,d}
      \cdot \frac{D_{c,\,v_{\text{home}}}}{\mathrm{FB}_{v_{\text{home}}}}
  \label{eq:omegaB}
\end{equation}

\textbf{Important:} $\Omega^B$ has the same algebraic form as $\Omega^A$
but a different interpretation. In $\Omega^A$, both $\lambda$ and
$\kappa$ refer to the \emph{same city} ($v$); the question is how many
of the arrivals at $v$ contact co-nationals in $v$. In $\Omega^B$,
$\lambda$ refers to the \emph{match venue} ($v_{\text{match}}$) and
$\kappa$ refers to the \emph{hub city} ($v_{\text{home}}$); these are
different places. $\Omega^B$ is therefore better interpreted as a
\emph{relative seeding index} --- it ranks hub cities by their exposure
potential, proportional to both the infectious pressure at the venue
and the density of the diaspora in the hub --- rather than a strict
Poisson rate for secondary cases.

The hub-city aggregate across all source countries and match venues is:
\[
  \Omega^B_{v_{\text{home}},\,d}
    = \sum_c \sum_{v_{\text{match}} \neq v_{\text{home}}}
      \Omega^B_{c,\,v_{\text{match}},\,v_{\text{home}},\,d}
\]

% ============================================================
\section*{5. Summary: what each quantity means}
% ============================================================

\begin{center}
\small
\begin{tabular}{p{2.8cm} p{4.2cm} p{5.8cm}}
\toprule
\textbf{Quantity} & \textbf{Definition} & \textbf{Plain interpretation} \\
\midrule
$\lambda_{c,v,d}$ &
  $N_{c,v} \cdot \rho_d \cdot (C_{c,d}/P_c) \cdot p_d$ &
  Expected infectious travelers from $c$ arriving at $v$. \\[6pt]
$P(\geq 1)_{v,d}$ &
  $1 - e^{-\Lambda_{v,d}}$ &
  Probability $\geq 1$ case arrives at city $v$. \\[6pt]
$\kappa_{c,v}$ &
  $D_{c,v} / \mathrm{FB}_{v}$ &
  Share of city $v$'s immigrants born in $c$. \\[6pt]
$\Omega^A_{c,v,d}$ &
  $\lambda_{c,v,d} \cdot \kappa_{c,v}$ &
  Expected cases \emph{entering co-national network} in $v$.
  Valid Poisson rate (thinned from $\lambda$). \\[6pt]
$P^A(\geq 1)_{v,d}$ &
  $1 - e^{-\Omega^A_{v,d}}$ &
  Probability $\geq 1$ case enters a co-national network in $v$.
  \textbf{Directly comparable to $P(\geq 1)_{v,d}$.} \\[6pt]
$\Omega^B_{c,v_m,v_h,d}$ &
  $\lambda_{c,v_m,d} \cdot \kappa_{c,v_h}$ &
  Relative seeding index for hub city $v_h$ from match at $v_m$.
  Ranks cities; not a probability. \\
\bottomrule
\end{tabular}
\end{center}

\bigskip
\noindent\textbf{Key constraint.} $\Omega^A_{c,v,d} \leq \lambda_{c,v,d}$
always, because $\kappa_{c,v} \leq 1$. The diaspora framework does not
create additional cases --- it partitions the existing importation risk
into the fraction that enters co-national networks ($\Omega^A$) versus
the fraction that disperses into the general population
($\lambda - \Omega^A$). The sum $\Omega^A + (\lambda - \Omega^A) = \lambda$
is exact.

\end{document}
